%% HAND-WRITTEN fragment -- NOT generated by maestro2tex.
%% Source of every number: the archived headless run
%%   ~/chips/castalia/simruns/cv_bc_37C/   (electrode Temp=310.15, matching temp=37)
%%     summary_ps_cv.txt                                  -> i_peak
%%     castalia/tb_anatop_pixel/maestro/results/maestro/Ocean.0/1/ps_cv/
%%       netlist/input.scs                                -> topology, tran line,
%%                                                           temp, ECS params
%%       psf/variables_file                               -> design variables
%% Predicted peak: Randles-Sevcik at the ECS parameters plus Cdl*dv/dt.
%% Sign convention measured on this run; the TB-07 formula in
%% ~/vestarv/docs/afe_rev2_simplan.html is inverted.

\paragraph{First measurement with the rev-2 channel in the loop.} The
cyclic-voltammetry test is the
potentiostat bench of Section~\ref{ss:pixel} with the cell master swapped from
the linear cell to the nonlinear electrode cell and nothing else changed:
the class-AB amplifier of Section~\ref{ss:ampab} as both A1 and A2, the local
bias generator of Section~\ref{ss:biasgen-local}, an ideal
\(R_f = \SI{35}{\kilo\ohm}\) rather than the programmable resistor of
Section~\ref{ss:tiarprog}, and \(C_{\mathrm{dl}} = \SI{0.2}{\micro\farad}\),
\(R_u = R_{\mathrm{ce}} = \SI{1}{\kilo\ohm}\) around an electrode at
\(E^{0\prime} = \SI{0.3}{\volt}\), \SI{1}{\milli\mole\per\litre},
\(k^0 = \SI{0.1}{\centi\metre\per\second}\). The commanded potential is a
\SIrange{0.75}{1.75}{\volt} triangle of \SI{0.5}{\second} per ramp, i.e.
\SI{2}{\volt\per\second}; transient to \SI{1.2}{\second} at
\texttt{errpreset=conservative}, typical process point, \SI{37}{\celsius} for
silicon and electrode alike (the model's \texttt{Temp} parameter is set to
\SI{310.15}{\kelvin} to match the simulator temperature).

The measured peak is \SI{8.666}{\micro\ampere} against
\SI{8.74}{\micro\ampere} predicted --- \SI{8.34}{\micro\ampere} of
Randles--\v{S}ev\v{c}\'ik faradaic current evaluated at \SI{310.15}{\kelvin}
plus \SI{0.40}{\micro\ampere} of \(C_{\mathrm{dl}}\,\mathrm{d}v/\mathrm{d}t\)
--- a \SI{-0.8}{\percent} discrepancy, so the nonlinearity survives the loop
rather than being reshaped by it. Current is read from the output voltage as
\(i_{\mathrm{we}} = (v_{\mathrm{out}} - V_{\mathrm{CM}})/R_f\): A2 holds WE at
\(V_{\mathrm{CM}}\) and the whole electrode current returns through \(R_f\), so
\(v_{\mathrm{out}} = V_{\mathrm{CM}} + i_{\mathrm{we}}R_f\). The peak therefore
moves the output \SI{0.30}{\volt} from \(V_{\mathrm{CM}} = \SI{1.25}{\volt}\),
\SI{38}{\percent} of the \(\pm\SI{0.80}{\volt}\) transimpedance window at this
gain, so the sweep is measured unclipped. The four wound-panel analytes run as
Maestro corners on the same test (each corner overrides
\(E^{0\prime}\), \(k^0\), \(D\), \(n\) and the bulk concentrations);
Figure~\ref{fig:electrode-panel} shows the resulting voltammograms.

\paragraph{What these numbers do not cover.} Five limits, each of which changes
how a figure above must be read.

\begin{itemize}[nosep]
  \item The in-loop results are nominal transients: one process point, one
    temperature, no process-corner set and no Monte Carlo. The analyte corners
    of Figure~\ref{fig:electrode-panel} vary only the electrode parameters and
    all run at the nominal process point.
  \item A Verilog-A element carries no mismatch statistics. The electrode model
    contributes nothing to a Monte Carlo run and must never be the thing one
    measures; every Monte Carlo \(\sigma\) in this chapter holds the electrode
    exactly nominal and reports circuit variation alone.
  \item The standalone validation of Table~\ref{tab:electrode-peaks} runs at
    \SI{25}{\celsius} throughout; the in-loop run above is the only result at
    \SI{37}{\celsius}. \texttt{Temp} is a plain model parameter, not tied to
    the simulator temperature --- any new bench must set it explicitly or the
    interface kinetics silently revert to \SI{25}{\celsius}. \(D\) is held at
    its \SI{25}{\celsius} value in both cases; the \SIrange{25}{37}{\celsius}
    increase in a real aqueous \(D\) (\(\sim\)30\,\%) is not modelled.
  \item Randles--\v{S}ev\v{c}\'ik assumes semi-infinite diffusion; the model
    pins bulk concentration at \(\delta = \SI{1.79}{\milli\metre}\). The
    difference is negligible while \(\delta \gg \sqrt{Dt_{\mathrm{scan}}}\),
    which holds by \(16\times\) at \SI{50}{\milli\volt\per\second}
    (\(\sqrt{Dt} = \SI{0.11}{\milli\metre}\)) and by far more at
    \SI{2}{\volt\per\second}. It does not hold for a thin unstirred film: a
    \SI{93}{\micro\metre} variant of the model exists for that case and
    attenuates the reverse peak to \(0.65\) of the forward one. Quote \(\delta\)
    with any peak taken from this model.
  \item Peak position and height identify and quantify an analyte; \(k^0\)
    only shapes the peak, and it is the least-supported parameter of the set.
    No heterogeneous rate constant in \si{\centi\metre\per\second} is published
    on carbon for any of the four wound markers the rev-2 panel targets --- the
    values in use are fitted to reproduce a measured \(\Delta E_p\) or
    \(\mathrm{d}E_p/\mathrm{d}\log v\). That is a gap in the literature, not in
    the model.
\end{itemize}
